\lampiransection{Instrumen dan Rekapitulasi Angket Tertutup IFAS dan EFAS}
\label{app:instrumen-rekapitulasi-angket-ifas-efas}

Lampiran ini memuat lembar fisik instrumen angket tertutup IFAS (\textit{Internal Factor Analysis Summary}) dan EFAS (\textit{External Factor Analysis Summary}) serta tabel rekapitulasi skor penilaian responden yang digunakan untuk menentukan bobot relatif, rating rata-rata, dan skor tertimbang dalam analisis SWOT.

\subsection{Lembar Fisik Instrumen Angket Tertutup IFAS dan EFAS}
Instrumen angket tertutup ini disusun untuk memperoleh penilaian kuantitatif mengenai tingkat kepentingan (bobot) dan tingkat kinerja/respons (rating) terhadap faktor internal dan eksternal Rumah Makan Nasi Gerilya. Petunjuk pengisian dan format instrumen disajikan pada Tabel~\ref{tab:instrumen-angket-ifas} dan Tabel~\ref{tab:instrumen-angket-efas}.

\paragraph{Petunjuk Pengisian Angket:}
\begin{enumerate}
    \item \textbf{Tingkat Kepentingan (Bobot):} Berikan nilai 1 sampai 5 untuk menunjukkan seberapa penting faktor ini mempengaruhi keberhasilan strategi pemasaran pada GrabFood (1 = Sangat Rendah, 2 = Rendah, 3 = Cukup, 4 = Tinggi, 5 = Sangat Tinggi).
    \item \textbf{Tingkat Kinerja / Respons (Rating):} 
    \begin{itemize}
        \item Untuk Faktor Internal: Berikan nilai 1 sampai 4 untuk menunjukkan kondisi usaha (1 = Kelemahan Utama, 2 = Kelemahan Pendukung, 3 = Kekuatan Pendukung, 4 = Kekuatan Utama).
        \item Untuk Faktor Eksternal: Berikan nilai 1 sampai 4 untuk menunjukkan respons usaha terhadap lingkungan (1 = Respons Sangat Lemah, 2 = Respons Cukup, 3 = Respons Baik, 4 = Respons Sangat Baik).
    \end{itemize}
\end{enumerate}

\begingroup
\scriptsize
\begin{xltabular}{\textwidth}{|C{0.8cm}|Y|C{2.6cm}|C{2.3cm}|}
\caption{Format Instrumen Angket Tertutup Penilaian Faktor Internal (IFAS)}
\label{tab:instrumen-angket-ifas}\\
\toprule
\textbf{No.} & \textbf{Pernyataan Faktor Internal} & \textbf{Skala Kepentingan (1--5)} & \textbf{Skala Kinerja (1--4)} \\ \hline
\midrule
\endfirsthead
\continuedtabletitle{4}
\toprule
\textbf{No.} & \textbf{Pernyataan Faktor Internal} & \textbf{Skala Kepentingan (1--5)} & \textbf{Skala Kinerja (1--4)} \\ \hline
\midrule
\endhead
\continuedtablefooter{4}
\endfoot
\bottomrule
\endlastfoot
\multicolumn{4}{|c|}{\cellcolor{black!10}\textbf{Kekuatan (\textit{Strengths})}} \\ \hline
1 & Visi jelas mengenai dapur sentral dan komitmen menjaga mutu rasa & [1][2][3][4][5] & [1][2][3][4] \\ \hline
2 & Cita rasa lauk kuat dan khas menjadi daya tarik utama produk & [1][2][3][4][5] & [1][2][3][4] \\ \hline
3 & Porsi makanan besar serta dukungan armada cold storage & [1][2][3][4][5] & [1][2][3][4] \\ \hline
4 & Pembagian struktur kerja per divisi serta respons cepat terhadap keluhan & [1][2][3][4][5] & [1][2][3][4] \\ \hline
5 & Reputasi digital toko yang baik (rating 4,7 di GrabFood) & [1][2][3][4][5] & [1][2][3][4] \\ \hline
6 & Pelaksanaan \textit{checkpoint} pesanan dan adopsi sistem digital & [1][2][3][4][5] & [1][2][3][4] \\ \hline
\multicolumn{4}{|c|}{\cellcolor{black!10}\textbf{Kelemahan (\textit{Weaknesses})}} \\ \hline
1 & Risiko kesalahan/kerawanan \textit{packing} item terpisah saat ramai & [1][2][3][4][5] & [1][2][3][4] \\ \hline
2 & Penumpukan antrean tiga jalur (dine in, take away, ojol) saat \textit{peak hour} & [1][2][3][4][5] & [1][2][3][4] \\ \hline
3 & Pembaruan status ketersediaan stok menu di aplikasi belum otomatis & [1][2][3][4][5] & [1][2][3][4] \\ \hline
4 & Persepsi harga GrabFood yang dianggap relatif mahal & [1][2][3][4][5] & [1][2][3][4] \\ \hline
5 & Deskripsi rincian isi paket menu pada aplikasi belum cukup tegas & [1][2][3][4][5] & [1][2][3][4] \\ \hline
6 & Keterbatasan margin keuntungan yang belum mencapai BEP ekspansi & [1][2][3][4][5] & [1][2][3][4] \\ \hline
\end{xltabular}
\tablesource[\textwidth]{Lembar instrumen angket tertutup IFAS penelitian (2026).}
\endgroup

\begingroup
\scriptsize
\begin{xltabular}{\textwidth}{|C{0.8cm}|Y|C{2.6cm}|C{2.3cm}|}
\caption{Format Instrumen Angket Tertutup Penilaian Faktor Eksternal (EFAS)}
\label{tab:instrumen-angket-efas}\\
\toprule
\textbf{No.} & \textbf{Pernyataan Faktor Eksternal} & \textbf{Skala Kepentingan (1--5)} & \textbf{Skala Respons (1--4)} \\ \hline
\midrule
\endfirsthead
\continuedtabletitle{4}
\toprule
\textbf{No.} & \textbf{Pernyataan Faktor Eksternal} & \textbf{Skala Kepentingan (1--5)} & \textbf{Skala Respons (1--4)} \\ \hline
\midrule
\endhead
\continuedtablefooter{4}
\endfoot
\bottomrule
\endlastfoot
\multicolumn{4}{|c|}{\cellcolor{black!10}\textbf{Peluang (\textit{Opportunities})}} \\ \hline
1 & Tingginya permintaan pesan-antar makanan praktis saat jam makan siang & [1][2][3][4][5] & [1][2][3][4] \\ \hline
2 & Pemanfaatan promo platform GrabFood untuk akuisisi pelanggan baru & [1][2][3][4][5] & [1][2][3][4] \\ \hline
3 & Ketersediaan pasokan bahan segar dan dukungan armada cold storage & [1][2][3][4][5] & [1][2][3][4] \\ \hline
4 & Peluang diferensiasi layanan melalui ketelitian \textit{packing} dibanding kompetitor & [1][2][3][4][5] & [1][2][3][4] \\ \hline
5 & Perubahan tren gaya hidup pemesanan digital dan program loyalitas pelanggan & [1][2][3][4][5] & [1][2][3][4] \\ \hline
6 & Integrasi sistem POS digital dan metode pembayaran QRIS & [1][2][3][4][5] & [1][2][3][4] \\ \hline
\multicolumn{4}{|c|}{\cellcolor{black!10}\textbf{Ancaman (\textit{Threats})}} \\ \hline
1 & Tingginya reputasi digital dan persaingan ketat restoran pesaing sejenis & [1][2][3][4][5] & [1][2][3][4] \\ \hline
2 & Kemudahan pelanggan membandingkan harga antar \textit{merchant} pada platform & [1][2][3][4][5] & [1][2][3][4] \\ \hline
3 & Risiko ulasan negatif pelanggan akibat kesalahan pengiriman atau item kurang & [1][2][3][4][5] & [1][2][3][4] \\ \hline
4 & Inflasi harga bahan baku serta skema komisi platform yang menekan margin & [1][2][3][4][5] & [1][2][3][4] \\ \hline
5 & Regulasi komisi platform serta kebijakan perpajakan daerah (PB1) & [1][2][3][4][5] & [1][2][3][4] \\ \hline
\end{xltabular}
\tablesource[\textwidth]{Lembar instrumen angket tertutup EFAS penelitian (2026).}
\endgroup

\clearpage
\subsection{Tabel Rekapitulasi Skor Angket Responden}
Angket tertutup IFAS khusus diisi oleh 3 responden internal yang memahami operasional dan kondisi internal usaha, yaitu Richard (Pemilik Usaha), Ririn (Kasir/Checker), dan Bella (Leader Dapur). Sementara itu, angket tertutup EFAS diisi oleh 3 responden internal tersebut serta ditambahkan partisipasi 3 responden pelanggan GrabFood (Alip, Reyza, dan Ojan) untuk menilai lingkungan eksternal. Rekapitulasi jawaban responden untuk bobot kepentingan (skala 1--5) dan rating kinerja/respons (skala 1--4) disajikan pada Tabel~\ref{tab:rekapitulasi-angket-responden-ifas} dan Tabel~\ref{tab:rekapitulasi-angket-responden-efas}.

\begin{landscape}
\begingroup
\scriptsize
\setlength{\tabcolsep}{4pt}
\begin{xltabular}{\linewidth}{|C{0.8cm}|Y|C{1.2cm}|C{1.2cm}|C{1.2cm}|C{1.6cm}|C{1.2cm}|C{1.2cm}|C{1.2cm}|C{1.6cm}|}
\caption{Rekapitulasi Penilaian Responden IFAS (Bobot dan Rating)}
\label{tab:rekapitulasi-angket-responden-ifas}\\
\toprule
 &  & \multicolumn{4}{c|}{\textbf{Tingkat Kepentingan / Bobot (1--5)}} & \multicolumn{4}{c|}{\textbf{Tingkat Kinerja / Rating (1--4)}} \\ \cline{3-10}
\textbf{No.} & \textbf{Faktor Internal} & \textbf{R1} & \textbf{R2} & \textbf{R3} & \textbf{Rerata} & \textbf{R1} & \textbf{R2} & \textbf{R3} & \textbf{Rerata} \\ \hline
\midrule
\endfirsthead
\continuedtabletitle{10}
\toprule
 &  & \multicolumn{4}{c|}{\textbf{Tingkat Kepentingan / Bobot (1--5)}} & \multicolumn{4}{c|}{\textbf{Tingkat Kinerja / Rating (1--4)}} \\ \cline{3-10}
\textbf{No.} & \textbf{Faktor Internal} & \textbf{R1} & \textbf{R2} & \textbf{R3} & \textbf{Rerata} & \textbf{R1} & \textbf{R2} & \textbf{R3} & \textbf{Rerata} \\ \hline
\midrule
\endhead
\continuedtablefooter{10}
\endfoot
\bottomrule
\endlastfoot
\multicolumn{10}{|c|}{\cellcolor{black!10}\textbf{Kekuatan (\textit{Strengths})}} \\ \hline
1 & Visi dapur sentral \& komitmen mutu rasa & 5 & 5 & 5 & 5,00 & 4 & 4 & 4 & 4,00 \\ \hline
2 & Cita rasa lauk kuat \& khas & 5 & 5 & 5 & 5,00 & 4 & 4 & 4 & 4,00 \\ \hline
3 & Porsi besar \& armada cold storage & 4 & 4 & 4 & 4,00 & 4 & 4 & 4 & 4,00 \\ \hline
4 & Struktur kerja per divisi \& respons keluhan & 4 & 4 & 4 & 4,00 & 3 & 3 & 3 & 3,00 \\ \hline
5 & Reputasi digital kuat (Rating 4,7) & 4 & 4 & 4 & 4,00 & 4 & 3 & 3 & 3,33 \\ \hline
6 & Checkpoint pesanan \& adopsi digital & 4 & 4 & 4 & 4,00 & 3 & 3 & 3 & 3,00 \\ \hline
\multicolumn{10}{|c|}{\cellcolor{black!10}\textbf{Kelemahan (\textit{Weaknesses})}} \\ \hline
1 & Risiko kesalahan packing item terpisah & 5 & 5 & 5 & 5,00 & 1 & 1 & 1 & 1,00 \\ \hline
2 & Penumpukan antrean 3 jalur peak hour & 4 & 4 & 4 & 4,00 & 1 & 1 & 1 & 1,00 \\ \hline
3 & Status stok menu aplikasi belum otomatis & 4 & 4 & 4 & 4,00 & 2 & 2 & 2 & 2,00 \\ \hline
4 & Harga GrabFood dipersepsikan mahal & 4 & 4 & 3 & 3,67 & 2 & 2 & 2 & 2,00 \\ \hline
5 & Deskripsi rincian paket belum tegas & 3 & 3 & 3 & 3,00 & 2 & 2 & 2 & 2,00 \\ \hline
6 & Margin belum mencapai BEP ekspansi & 3 & 4 & 3 & 3,33 & 2 & 2 & 2 & 2,00 \\ \hline
\multicolumn{2}{|c|}{\textbf{Total Rata-rata Kepentingan}} & \multicolumn{4}{c|}{\textbf{49,00}} & \multicolumn{4}{c|}{--} \\ \hline
\end{xltabular}
\tablesource[\linewidth]{Keterangan Responden Internal IFAS: R1 = Richard (Pemilik), R2 = Ririn (Kasir/Checker), R3 = Bella (Leader Dapur).}
\endgroup
\end{landscape}

\begin{landscape}
\begingroup
\scriptsize
\setlength{\tabcolsep}{3pt}
\begin{xltabular}{\linewidth}{|C{0.7cm}|Z{5.2cm}|C{0.7cm}|C{0.7cm}|C{0.7cm}|C{0.7cm}|C{0.7cm}|C{0.7cm}|C{1.1cm}|C{0.7cm}|C{0.7cm}|C{0.7cm}|C{0.7cm}|C{0.7cm}|C{0.7cm}|C{1.1cm}|}
\caption{Rekapitulasi Penilaian Responden EFAS (Bobot dan Rating)}
\label{tab:rekapitulasi-angket-responden-efas}\\
\toprule
 &  & \multicolumn{7}{c|}{\textbf{Tingkat Kepentingan / Bobot (1--5)}} & \multicolumn{7}{c|}{\textbf{Tingkat Respons (1--4)}} \\ \cline{3-16}
\textbf{No.} & \textbf{Faktor Eksternal} & \textbf{R1} & \textbf{R2} & \textbf{R3} & \textbf{R4} & \textbf{R5} & \textbf{R6} & \textbf{Rerata} & \textbf{R1} & \textbf{R2} & \textbf{R3} & \textbf{R4} & \textbf{R5} & \textbf{R6} & \textbf{Rerata} \\ \hline
\midrule
\endfirsthead
\continuedtabletitle{16}
\toprule
 &  & \multicolumn{7}{c|}{\textbf{Tingkat Kepentingan / Bobot (1--5)}} & \multicolumn{7}{c|}{\textbf{Tingkat Respons (1--4)}} \\ \cline{3-16}
\textbf{No.} & \textbf{Faktor Eksternal} & \textbf{R1} & \textbf{R2} & \textbf{R3} & \textbf{R4} & \textbf{R5} & \textbf{R6} & \textbf{Rerata} & \textbf{R1} & \textbf{R2} & \textbf{R3} & \textbf{R4} & \textbf{R5} & \textbf{R6} & \textbf{Rerata} \\ \hline
\midrule
\endhead
\continuedtablefooter{16}
\endfoot
\bottomrule
\endlastfoot
\multicolumn{16}{|c|}{\cellcolor{black!10}\textbf{Peluang (\textit{Opportunities})}} \\ \hline
1 & Permintaan makan siang praktis & 5 & 5 & 5 & 5 & 5 & 5 & 5,00 & 3 & 3 & 3 & 3 & 3 & 3 & 3,00 \\ \hline
2 & Promo GrabFood akuisisi baru & 5 & 5 & 5 & 5 & 5 & 5 & 5,00 & 3 & 3 & 3 & 3 & 3 & 3 & 3,00 \\ \hline
3 & Pasokan bahan segar \& cold storage & 4 & 4 & 4 & 4 & 4 & 4 & 4,00 & 4 & 3 & 3 & 3 & 3 & 3 & 3,17 \\ \hline
4 & Kelemahan packing kompetitor & 4 & 4 & 4 & 4 & 4 & 4 & 4,00 & 3 & 3 & 3 & 3 & 3 & 3 & 3,00 \\ \hline
5 & Tren pemesanan digital \& loyalitas & 4 & 4 & 4 & 4 & 4 & 4 & 4,00 & 3 & 3 & 3 & 3 & 3 & 3 & 3,00 \\ \hline
6 & Integrasi POS digital \& QRIS & 4 & 4 & 4 & 4 & 4 & 4 & 4,00 & 3 & 3 & 3 & 3 & 3 & 3 & 3,00 \\ \hline
\multicolumn{16}{|c|}{\cellcolor{black!10}\textbf{Ancaman (\textit{Threats})}} \\ \hline
1 & Persaingan \& reputasi pesaing & 5 & 5 & 5 & 5 & 5 & 5 & 5,00 & 2 & 2 & 2 & 2 & 2 & 2 & 2,00 \\ \hline
2 & Pembandingan harga antar-merchant & 4 & 4 & 4 & 4 & 4 & 4 & 4,00 & 2 & 2 & 2 & 2 & 2 & 2 & 2,00 \\ \hline
3 & Risiko ulasan negatif pengiriman & 4 & 4 & 4 & 4 & 4 & 4 & 4,00 & 2 & 2 & 2 & 2 & 2 & 2 & 2,00 \\ \hline
4 & Inflasi bahan baku \& komisi & 4 & 4 & 4 & 4 & 4 & 4 & 4,00 & 2 & 2 & 2 & 2 & 2 & 2 & 2,00 \\ \hline
5 & Regulasi komisi \& pajak PB1 & 3 & 3 & 3 & 3 & 3 & 3 & 3,00 & 2 & 2 & 2 & 2 & 2 & 2 & 2,00 \\ \hline
\multicolumn{2}{|c|}{\textbf{Total Rata-rata Kepentingan}} & \multicolumn{7}{c|}{\textbf{46,00}} & \multicolumn{7}{c|}{--} \\ \hline
\end{xltabular}
\tablesource[\linewidth]{Keterangan Responden: R1 = Richard (Pemilik), R2 = Ririn (Kasir/Checker), R3 = Bella (Leader Dapur), R4 = Alip (Pelanggan), R5 = Reyza (Pelanggan), R6 = Ojan (Pelanggan).}
\endgroup
\end{landscape}

\subsection{Tabel Perhitungan Bobot Relatif, Rating, dan Skor Tertimbang}
Tabel~\ref{tab:perhitungan-bobot-rating-skor-ifas-efas} menyajikan proses konversi dari rata-rata nilai kepentingan responden menjadi bobot relatif (normalisasi total 1,00) serta perhitungan skor tertimbang akhir untuk IFAS dan EFAS yang digunakan di Bab IV.

\begingroup
\scriptsize
\begin{xltabular}{\textwidth}{|C{0.8cm}|Y|C{2.0cm}|C{2.0cm}|C{1.8cm}|C{2.2cm}|}
\caption{Tabel Perhitungan Bobot Relatif, Rating, dan Skor Tertimbang IFAS dan EFAS}
\label{tab:perhitungan-bobot-rating-skor-ifas-efas}\\
\toprule
\textbf{No.} & \textbf{Faktor Strategis} & \textbf{Rerata Kepentingan} & \textbf{Bobot Relatif} & \textbf{Rating Rerata} & \textbf{Skor Tertimbang} \\ \hline
\midrule
\endfirsthead
\continuedtabletitle{6}
\toprule
\textbf{No.} & \textbf{Faktor Strategis} & \textbf{Rerata Kepentingan} & \textbf{Bobot Relatif} & \textbf{Rating Rerata} & \textbf{Skor Tertimbang} \\ \hline
\midrule
\endhead
\continuedtablefooter{6}
\endfoot
\bottomrule
\endlastfoot
\multicolumn{6}{|c|}{\cellcolor{black!10}\textbf{Faktor Internal (IFAS)}} \\ \hline
S1 & Visi dapur sentral \& komitmen mutu rasa & 5,00 & 0,10 & 4,00 & 0,40 \\ \hline
S2 & Cita rasa lauk kuat \& khas & 5,00 & 0,10 & 4,00 & 0,40 \\ \hline
S3 & Porsi besar \& armada cold storage & 4,00 & 0,08 & 4,00 & 0,32 \\ \hline
S4 & Struktur kerja per divisi \& respons keluhan & 4,00 & 0,08 & 3,00 & 0,24 \\ \hline
S5 & Reputasi digital kuat (Rating 4,7) & 4,00 & 0,08 & 3,33 & 0,27 \\ \hline
S6 & Checkpoint pesanan \& adopsi digital & 4,00 & 0,08 & 3,00 & 0,25 \\ \hline
W1 & Risiko kesalahan packing item terpisah & 5,00 & 0,10 & 1,00 & 0,10 \\ \hline
W2 & Penumpukan antrean 3 jalur peak hour & 4,00 & 0,08 & 1,00 & 0,08 \\ \hline
W3 & Status stok menu aplikasi belum otomatis & 4,00 & 0,08 & 2,00 & 0,16 \\ \hline
W4 & Harga GrabFood dipersepsikan mahal & 3,67 & 0,08 & 2,00 & 0,15 \\ \hline
W5 & Deskripsi rincian paket belum tegas & 3,00 & 0,06 & 2,00 & 0,12 \\ \hline
W6 & Margin belum mencapai BEP ekspansi & 3,33 & 0,08 & 2,00 & 0,13 \\ \hline
\textbf{Total} & \textbf{IFAS} & \textbf{49,00} & \textbf{1,00} & -- & \textbf{2,62} \\ \hline
\multicolumn{6}{|c|}{\cellcolor{black!10}\textbf{Faktor Eksternal (EFAS)}} \\ \hline
O1 & Permintaan makan siang praktis & 5,00 & 0,11 & 3,00 & 0,32 \\ \hline
O2 & Promo GrabFood akuisisi baru & 5,00 & 0,11 & 3,00 & 0,32 \\ \hline
O3 & Pasokan bahan segar \& cold storage & 4,00 & 0,09 & 3,17 & 0,26 \\ \hline
O4 & Kelemahan packing kompetitor & 4,00 & 0,09 & 3,00 & 0,24 \\ \hline
O5 & Tren pemesanan digital \& loyalitas & 4,00 & 0,09 & 3,00 & 0,24 \\ \hline
O6 & Integrasi POS digital \& QRIS & 4,00 & 0,09 & 3,00 & 0,24 \\ \hline
T1 & Persaingan \& reputasi pesaing & 5,00 & 0,11 & 2,00 & 0,22 \\ \hline
T2 & Pembandingan harga antar-merchant & 4,00 & 0,09 & 2,00 & 0,18 \\ \hline
T3 & Risiko ulasan negatif pengiriman & 4,00 & 0,09 & 2,00 & 0,18 \\ \hline
T4 & Inflasi bahan baku \& komisi & 4,00 & 0,09 & 2,00 & 0,18 \\ \hline
T5 & Regulasi komisi \& pajak PB1 & 3,00 & 0,06 & 2,00 & 0,16 \\ \hline
\textbf{Total} & \textbf{EFAS} & \textbf{46,00} & \textbf{1,00} & -- & \textbf{2,54} \\ \hline
\end{xltabular}
\tablesource[\textwidth]{Diolah peneliti berdasarkan hasil rekapitulasi angket tertutup responden internal (IFAS) dan responden internal-pelanggan (EFAS) (2026).}
\endgroup
